
\section{Experiments}
\label{sec:experiments}


\begin{table}[t]
    \centering
    \resizebox{0.98\linewidth}{!}{
        \begin{tabular}{lcccc}
            \toprule
            & \tf{Fluency} & \tf{Correct.} & \multicolumn{2}{c}{\tf{Citation} }\\
            \cmidrule(lr){2-2} \cmidrule(lr){3-3} \cmidrule(lr){4-5}
            & (MAUVE) & (EM Rec.) & Rec. & Prec. \\
            \midrule
            \headercolor
            \multicolumn{5}{c}{\textbf{ChatGPT}} \\
            \tableskip
            \vani{} (5-psg) & 66.6 & 40.4  & {73.6}  & 72.5\\% {63.0} \\
            ~~w/ \rerank{} &  77.0 & 40.2 & \tf{84.8} & \textbf{81.6} \\% \tf{73.3} \\
            \summ{} (10-psg) & 70.0  & \tf{43.3}  &68.9	& 61.8\\ %
            ~~w/ \interact{}  & 69.0 &	39.1 & 73.4& 66.5\\ %
            \snippet{} (10-psg) & 69.8&	41.4& 65.3&	57.4\\ %
            \search{} & 58.7 & 32.4 & 58.3 & 58.2	\\ %
            \close{} & 52.7 &	38.3 & 26.7& 26.7 \\

            \tableskip
            \headercolor
             \multicolumn{5}{c}{\textbf{GPT-4} (\vani{} prompting)}  \\
            \tableskip
            GPT-4 (5-psg) & 67.1 & 41.3 & 68.5 & 75.6 \\
            GPT-4 (20-psg) & 64.9 & 44.4 & 73.0 & 76.5 \\

            \tableskip
            \headercolor
            \multicolumn{5}{c}{\textbf{LLaMA} (\vani{} prompting)} \\
            \tableskip
            LLaMA-13B (3-psg) & 68.4 & 26.9 & 10.6 & 15.4 \\
            Vicuna-13B (3-psg) & 82.6 & 31.9 & 51.1 & 50.1 \\ 
 
            Chat-13B (5-psg) & 72.4 & 35.2 & 38.4 & 39.4 \\
            Chat-70B (5-psg) & 88.3 & 41.5 & 62.9 & 61.3 \\
            \bottomrule
        \end{tabular}
    }
    \caption{
        Experiments on ASQA.
        For \close{}, we use \posthoc{} to get citations.
        $k$-psg: putting top-$k$ passages from the retrieval results into the context.
        Chat-13B and Chat-70B refer to LLaMA-2-Chat.
    }
    \label{tab:asqa_result}
    \vspace{-10pt}
\end{table}


\begin{table}[t]
    \centering
    \small
    \resizebox{0.98\linewidth}{!}{
        \begin{tabular}{lcccc}
            \toprule
             & \multicolumn{2}{c}{\tf{Correctness}} & \multicolumn{2}{c}{\tf{Citation} }\\
             \cmidrule(lr){2-3} \cmidrule(lr){4-5}
            &  Rec.-5 & Prec. & Rec. & Prec. \\
            \midrule
            \headercolor
            \multicolumn{5}{c}{\textbf{ChatGPT}} \\
            \tableskip
            \vani{} (5-psg) & 20.8 &  20.8 & 20.5&	20.9 \\
            ~~w/ \rerank{} & 22.8 & 21.4 & 21.2 & 21.4 \\ %
            \summ{} (10-psg) & 23.6& 21.2&   \tf{23.6}&	\tf{25.7} \\
            \snippet{} (10-psg)&24.5 &  21.5&  {22.9} &  {24.9} \\
            ~~w/ \interact{} & 21.9 &  \tf{23.0} & 21.9&23.4\\
            \search{} & 17.2& 20.4 & 14.9&	14.9 \\
            \close{} &  \tf{32.9} & 19.8 &  10.0 &   10.0\\
            \tableskip
            \headercolor
             \multicolumn{5}{c}{\textbf{GPT-4} (\vani{} prompting)}  \\
            \tableskip
            GPT-4 (5-psg) & 22.2 & 25.0 & 25.9 & 27.0 \\
            GPT-4 (20-psg) & 29.6 & 26.2 & 27.4 & 28.5 \\ 

            \tableskip
            \headercolor
            \multicolumn{5}{c}{\textbf{LLaMA} (\vani{} prompting)} \\
            \tableskip
            LLaMA-13B (3-psg) & 9.7 & 9.1 & 6.7 & 7.1 \\% 6.8 \\
            Vicuna-13B (5-psg) & 14.0 & 15.9 & 12.5 & 13.4 \\ %
            
            
            Chat-13B (5-psg) & 21.1 & 18.2 & 9.6 & 9.7 \\
            Chat-70B (5-psg) & 21.8 &18.4 & 15.1 &15.6 \\
            \bottomrule
        \end{tabular}
    }
    \caption{
        Experiments  on QAMPARI.
        ``Rec.-5'': we set the recall to be 100\% if the prediction includes at least 5 correct answers.
    }
    \label{tab:qampari_result}
    \vspace{-10pt}
\end{table}



\begin{table}[h]
    \centering
    \resizebox{0.98\linewidth}{!}{
        \begin{tabular}{lcccc}
            \toprule
            & \tf{Fluency} & \tf{Correct.} & \multicolumn{2}{c}{\tf{Citation} }\\
            \cmidrule(lr){2-2} \cmidrule(lr){3-3} \cmidrule(lr){4-5}
            & (MAUVE) & (Claim) & Rec. & Prec. \\
            \midrule
            \headercolor
            \multicolumn{5}{c}{\textbf{ChatGPT}} \\
            \tableskip
            \vani{} (5-psg) &57.2&	12.0 &	51.1&	{50.0} \\
            ~~w/ \rerank{} & 56.1 & 11.4 & \tf{69.3} & \textbf{67.8}\\ %
            \summ{} (10-psg) & 40.3 & 12.5 & 51.5 & 48.2 \\
            \snippet{} (10-psg) & 62.9 &	{14.3} & 50.4 & 45.0 \\ %
            ~~w/ \interact{} & 68.0 &	13.3&	47.8& 45.0\\	%
            \search{} & 49.7&	13.4&	45.6&  43.7\\	%
            \close{} & 32.6 &\tf{18.6}	&15.5	&15.5\\
            \tableskip
            \headercolor
             \multicolumn{5}{c}{\textbf{GPT-4} (\vani{} prompting)}  \\
            \tableskip
            GPT-4 (5-psg) & 38.4 & 14.2 & 44.0 & 50.1 \\
            GPT-4 (20-psg) & 41.5 & 18.3 & 48.5 & 53.4 \\

            \tableskip
            \headercolor
            \multicolumn{5}{c}{\textbf{LLaMA} (\vani{} prompting)} \\
            \tableskip
            LLaMA-13B (3-psg) & 50.0 & 3.9 & 3.1 & 5.3 \\%2.0 \\
            Vicuna-13B (3-psg)& 58.2 &  {10.0} & 15.6 & 19.6 \\% 12.6 \\
            
            Chat-13B (5-psg) & 34.7 & 13.4 &17.3 & 15.8 \\
            Chat-70B (5-psg) & 38.6 & 12.8 & 38.3 & 37.9 \\
            \bottomrule
        \end{tabular}
    }
    \caption{
        Experiments on ELI5. We use \textit{claim recall} for the correctness evaluation.
        Chat-13B and Chat-70B refer to LLaMA-2-Chat.
    }
    \label{tab:eli5_result}
    \vspace{-10pt}
\end{table}


We describe experiment details in \S\ref{app_sec:imp_details}. 
We use ChatGPT (\ttt{gpt-3.5-turbo-0301}) with a 4K context window for most main experiments and ablations. 
\revise{We also report results with ChatGPT-16K (\ttt{gpt- 3.5-turbo-16k-0613}) and GPT-4 (\ttt{gpt-4-0613}; 8K context window).}
For open-source models, we test LLaMA~\citep{touvron2023llama} and 
its instruction-tuned versions, including 
Alpaca~\citep{alpaca}, Vicuna~\citep{vicuna2023}, and Oasst~\citep{kopf2023openassistant}.
They all have a 2K context window.
We use short instructions for LLaMA (Table~\ref{tab:inst_vanilla_light}) to save  context budget.
\revise{Additionally, we test LLaMA-2-Chat, which were also trained to follow instructions \cite{touvron2023llama2}. These models have a context window of 4K tokens, which allows for 5 passages per question.}



\subsection{Main Results}
We present the main results on three datasets in 
Table~\ref{tab:asqa_result},~\ref{tab:qampari_result}, and~\ref{tab:eli5_result} respectively (full results in \S\ref{app_sec:more_main}). %
We first note that all models achieve good fluency scores (except some models on ELI5 mainly due to their longer generations).
We summarize the main takeaways from the experiments below.

\paragraph{\tf{\textsc{\pvani{}}} achieves strong performance.}
Despite its simplicity,
\vani{} (putting retrieved passages in context)
achieves close-to-the-best performance among all prompting strategies. %



\paragraph{Using summaries or snippets improves correctness.}
We see a universal trend that 
\summ{} or \snippet{} 
improves  correctness,
though
on ASQA and ELI5, such an improvement comes at a cost of citation quality
due to the lossy compression. %
Combining \interact{} with \summ{}/\snippet{} does not bring improvement,
and we hypothesize that 
checking the full passages offers limited benefit and 
 current LLMs are not proficient in an interactive usage.






\paragraph{Retrieving text on the fly does not improve performance.}
All datasets show that 
\vani{} 
 outperforms
\search{} on citation quality (and on correctness for ASQA and ELI5).
By manually examining the examples,
we find that 
it is challenging to ask detailed questions without seeing any passages.
To improve \search{},
one may need to provide more context about the questions in advance or
encourage the model to call retrievers with more detailed and diverse queries.

\paragraph{{\textsc{\prerank{}}} boosts citation quality.}
We observe that \rerank{} 
leads to consistent improvement in citation quality (on ASQA and ELI5).
As the automatic scores may be biased in \rerank{},
we also conduct human evaluation~(\S\ref{sec:human})
and verify its effectiveness. %

\paragraph{\textsc{\pclose{}}+\textsc{\pposthoc{}} delivers strong correctness but poor citation quality.}
\close{} outperforms \vani{} in correctness on ELI5 and QAMPARI, and has only a 2\% gap on ASQA.
However, 
\close{} cannot provide any citation;
when combined with \posthoc{},
the citation quality remains inadequate. For instance, 
citation recall of \close{}+\posthoc{} is lower than \vani{} by 47\% on ASQA. 

To understand why \close{} achieves better correctness and why \posthoc{} cannot deliver satisfying citation quality,
we manually examine model outputs and find that:
(1) open-book models are easily distracted by irrelevant passages and generate responses with lower correctness, 
a phenomenon also observed by~\citet{shi2023large};
(2) \close{} often generates texts that are correct but not similar to any retrieved passages, 
making it difficult to match a citation post-hoc.

\revise{
\paragraph{GPT-4 brings limited improvement but is better at using long context.}
We evaluate GPT-4 with \vani{} and different numbers of passages (more results in \S\ref{app_sec:more_main}).
GPT-4 brings consistent (but limited) improvement on correctness, but often at a cost of citation quality. 
GPT-4 can also incorporate more passages due to its longer context window,
which boosts both correctness and citation quality. 
On the contrary, 
including more passages with ChatGPT-16K does not improve the results (Table \ref{tab:asqa_llm}), 
suggesting that processing more passages is non-trivial and GPT-4 is better at synthesizing information from its long context than ChatGPT.
}

\subsection{Comparison of Different LLMs}
\label{sec:different_llms}


\begin{figure*}[t]
    \centering
    \begin{minipage}[b]{0.66\textwidth}
    \centering
    \includegraphics[width=1.0\textwidth]{figs/ret.pdf}
    \end{minipage}
    \hfill
    \begin{minipage}[b]{0.33\textwidth}
    \centering

    \resizebox{1.0\textwidth}{!}{
        \begin{tabular}{lcccc}
            \toprule
            & \tf{Fluency} & \tf{Correct.} & \multicolumn{2}{c}{\tf{Citation} }\\
            \cmidrule(lr){2-2} \cmidrule(lr){3-3} \cmidrule(lr){4-5}
            & (MAUVE) & (EM) & Rec. & Prec. \\
            \midrule
            \headercolor
            \multicolumn{5}{c}{\textbf{ChatGPT} (\vani{})} \\
            \tableskip
            DPR (5-psg) &  61.8 &	36.1 & 65.0 &	65.6\\%56.7 \\
            GTR (1-psg) & 69.5&	38.4 & 56.0	&       64.0\\%56.0\\
            GTR (3-psg) & 66.6 & 39.6 & 72.8	&   \tf{73.9}\\%64.0 \\
            GTR (5-psg) & 66.6 & \tf{40.4}  & \tf{73.6} &     72.5\\%63.0 \\
            \bottomrule
        \end{tabular}
    }
    \vspace{17pt}
    \end{minipage}
    \vspace{-25pt}
    \caption{
Retrieval recall@$k$ on ASQA (\textit{EM recall}), QAMPARI (\textit{recall-5}), and ELI5 (\textit{claim recall}).
Retrieval recall serves as an upper bound for model performance,
and we compare them with two models' correctness results in the figure (dashed lines):
``Vanilla (5-psg)'' is ChatGPT \vani{} with top-5 passages in context;
``Oracle'' is the same model  except that it uses 5 gold passages (\S\ref{app_sec:retrieval}), whose recall matches Recall@100 on all three datasets. %
    }
    \vspace{-10pt}
    \label{fig:ret}
\end{figure*}


Table~\ref{tab:asqa_llm} compares 
different LLMs on ASQA using \vani{} (more results  in \S\ref{app_sec:more_main}).
Notably, instruction-tuned models (Vicuna-13B and LLaMA-2-Chat)
outperform the original LLaMA models in correctness and 
considerably enhance the citation quality.
We observe that 
while the original LLaMA models are able to %
copy facts from the context,
they struggle with accurately citing the sources or simply do not cite. 
\revise{Notably, the best open-source model, LLaMA-2-70B-Chat, achieves comparable correctness score as the OpenAI models, but still lags behind in citation quality. }





\subsection{Retrieval Analysis}


\begin{table}[t]
    \centering
    \resizebox{0.98\linewidth}{!}{
        \begin{tabular}{lcccc}
            \toprule
            & \tf{Fluency} & \tf{Correct.} & \multicolumn{2}{c}{\tf{Citation} }\\
            \cmidrule(lr){2-2} \cmidrule(lr){3-3} \cmidrule(lr){4-5}
            & (MAUVE) & (EM Rec.) & Rec. & Prec. \\
            \midrule
            \headercolor
            \multicolumn{5}{c}{\textbf{Open-source} (max \#tokens=2K-4K)} \\
            \tableskip
            LLaMA-13B (3-psg) & 68.4 & 26.9 & 10.6 & 15.4 \\
            Vicuna-13B (3-psg) & 82.6 & 31.9 & 51.1 & 50.1 \\ 
            Chat-13B (5-psg) & 72.4 & 35.2 & 38.4 & 39.4 \\
            Chat-70B (5-psg) & 88.3 & 41.5 & 62.9 & 61.3 \\
            \tableskip
            \headercolor
            \multicolumn{5}{c}{\textbf{ChatGPT} (max \#tokens=4K)} \\
            \tableskip
            ChatGPT (3-psg) & 66.6 & 39.6 & 72.8	& 73.9 \\% 64.0 \\
            ChatGPT (5-psg) & 66.6 & 40.4  & {73.6}  & 72.5 \\%{63.0} \\
            \tableskip
            \headercolor
            \multicolumn{5}{c}{\textbf{ChatGPT-16K} (max \#tokens=16K)} \\
            \tableskip
            ChatGPT (5-psg) & 60.3 & 36.1 & 76.2 & 76.5 \\
            ChatGPT (10-psg) & 56.3 & 36.7 & 75.3 & 75.0 \\ 
            ChatGPT (20-psg) & 56.7 & 36.1 & 73.7 & 73.5 \\

            \tableskip
            \headercolor
             \multicolumn{5}{c}{\textbf{GPT-4} (max \#tokens=8K)} \\
            \tableskip
            GPT-4 (5-psg) & 67.1 & 41.3 & 68.5 & 75.6 \\
            GPT-4 (10-psg) & 71.5 & 43.1 & 72.0 & 75.5 \\
            GPT-4 (20-psg) & 64.9 & 44.4 & 73.0 & 76.5 \\
            \bottomrule
 
        \end{tabular}
    }
    \caption{
        Comparison of different LLMs on ASQA (GTR+\vani{}). LLaMA-13B and Vicuna-13B have a context limit of 2,048 tokens, and thus can only use a short version of instructions and at most top-3 passages.
        Chat-13B and Chat-70B refer to LLaMA-2-Chat.
    }
    \label{tab:asqa_llm}
    \vspace{-10pt}
\end{table}


The retrieval results play a crucial role to the correctness and the citation quality.
Figure~\ref{fig:ret} presents 
the retrieval recall@$k$ %
with different datasets and retrievers.
As the number of passages increases, retrieval recall steadily improves.
Additionally, Figure~\ref{fig:ret} shows
the correctness performance of two models:
(1) ChatGPT \vani{} with top-5 passages (our primary baseline);
(2) an oracle version of the same model employing
5 gold passages (\S\ref{app_sec:retrieval}; the 5 gold passages match the retrieval recall@100).
Notably,
both models' correctness lags behind the corresponding retrieval recall (except for ELI5 top-5).
The discrepancy suggests that 
despite the presence of accurate answers in context,
LLMs struggle to utilize  them in their outputs.



We compare the impact of different retrievers and different numbers of passages to LLMs.
Figure~\ref{fig:ret} (right)
shows that
GTR outperforms DPR in both correctness and citation quality,
emphasizing  the importance of deploying better retrievers.
Contrary to the  retrieval recall trend  in Figure~\ref{fig:ret},
more passages in context do not yield substantial improvement for ChatGPT.
Specifically, correctness plateaus at top-1 passage and citation quality plateaus at top-3. 
\revise{GPT-4 (Table~\ref{tab:asqa_llm}) exhibits an increasing trend with more passages, 
but the improvement is not proportional to the retrieval performance.}
This  indicates the limited ability of LLMs in  utilizing multiple passages within context.

\subsection{Other Ablations}
\label{sec:other_ablation}

We provide additional ablations in \S\ref{app_sec:more_exp}. 
In summary, we find that 
(1) 
using comprehensive instructions
 enhances the citation quality of instruction-tuned models (\S\ref{sec:inst});
(2) 
including at least one demonstration improves the performance %
(\S\ref{sec:demo}); 
\revise{(3) fine-tuned models (FiD; \citealp{izacard-grave-2021-leveraging}) with \posthoc{} lag behind LLMs in both correctness and citation quality and fail to generalize %
(\S\ref{sec:fid}).}







